\subsection{EXPANSIONS OF A DETERMINANT}

Let I be a totally ordered finite indexing set. For every subset H of I let \(H'\) denote the complement I - H. Let \(X = (\xi_{ji})\) be a square matrix of type (I, I), which can be considered as the matrix of an endomorphism u of M = A \(^{I}\) with respect to the canonical basis \((e_{i})_{i \in I}\) of M. Let n = Card(I) and let H be a subset of I with \(q \leqslant n\) elements and K a subset of I with n - q elements; then we can write (no. 5, Proposition 10)

\[
\begin{array}{r l} (\bigwedge^ {q} (u)) (e _ {\mathrm{H}}) & = \sum_ {\mathbf {R}} \det (X _ {\mathbf {R}, \mathrm{H}}) e _ {\mathbf {R}} \\ (\bigwedge^ {n - q} (u)) (e _ {\mathbf {K}}) & = \sum_ {\mathbf {S}} \det (X _ {\mathbf {S}, \mathbf {K}}) e _ {\mathbf {S}} \end{array}
\]

where R (resp. S) runs through the set of subsets of I with q (resp. n - q) elements. It follows from § 7, no. 8, formulae (19) and (20) that

\[
e _ {\mathbf {R}} \wedge e _ {\mathbf {S}} = 0
\]

except when \(\mathbf{S} = \mathbf{R}'\) , whence the formula

\[
\left(\bigwedge^ {q} (u) \left(e _ {\mathrm{H}}\right)\right) \wedge \left(\bigwedge^ {n - q} (u) \left(e _ {\mathrm{K}}\right)\right) = \sum_ {\mathbf {R}} \rho_ {\mathbf {R}, \mathbf {R} ^ {\prime}} \det \left(X _ {\mathbf {R}, \mathbf {H}}\right) \det \left(X _ {\mathbf {R} ^ {\prime}, \mathbf {K}}\right) e _ {\mathbf {I}}\tag{20}
\]

where R runs through the set \(\mathfrak{F}_{q}(\mathbf{I})\) of subsets of I with q elements.

If we take \(\mathbf{K} = \mathbf{H}'\) , it follows from the definition of \(\bigwedge^{n}(u)\) ( \(\S 7\) , no. 2, formula (4)) and \(\S 7\) , no. 3, Corollary 1 to Proposition 5 that the right hand side of (20) is \(\rho_{\mathrm{H},\mathrm{H}'}\bigwedge^{n}(u)(e_{\mathrm{I}})\) . Hence (no. 1, formula (1) and \(\S 7\) , no. 2, formula (4))

\[
\det (X) = \rho_ {\mathrm{H}, \mathrm{H} ^ {\prime}} \sum_ {\mathbf {R} \in \mathfrak {F} _ {q} (\mathbf {I})} \rho_ {\mathrm{R}, \mathrm{R} ^ {\prime}} \det (X _ {\mathbf {R}, \mathrm{H}}) \det (X _ {\mathbf {R} ^ {\prime}, \mathrm{H} ^ {\prime}}).\tag{21}
\]

If on the other hand \(\mathbf{K} \neq \mathbf{H}'\) , then \(\mathbf{H} \cap \mathbf{K} \neq \varnothing\) ; as the left hand side of (20) is \(\pm \bigwedge^{n}(u)(e_{\mathbf{H}} \wedge e_{\mathbf{K}})\) , it is zero, whence

\[
\sum_ {\mathbf {R}} \rho_ {\mathbf {R}, \mathbf {R} ^ {\prime}} \det (X _ {\mathbf {R}, \mathbf {H}}) \det (X _ {\mathbf {R} ^ {\prime}, \mathbf {K}}) = 0 \quad \text { for } \mathbf {K} \neq \mathbf {H} ^ {\prime}.\tag{22}
\]

The right hand side of (21) is called the Laplace expansion of the determinant of the matrix \(X\) by the \(q\) columns whose indices belong to \(\mathbf{H}\) and the \(n - q\) columns whose indices belong to the complement \(\mathbf{H}'\) of \(\mathbf{H}\) . The minors \(\det(X_{\mathbf{R},\mathbf{H}})\) and \(\det(X_{\mathbf{R}',\mathbf{H}'})\) are sometimes called complementary.

An important simple case of the Laplace expansion is that where \(\mathbf{I} = [1, n]\) and \(q = 1\) , hence \(\mathrm{H} = \{i\}\) ; for every subset \(\mathbf{R} = \{j\}\) of \(\mathbf{I}\) with one element then \(\det X_{\mathbf{R},\mathbf{H}} = \xi_{ji}\) . The minor of \(\det X_{\mathbf{R}',\mathbf{H}'}\) is the determinant of the square matrix derived canonically (no. 5) from the matrix obtained by suppressing in \(X\) the row of index \(j\) and the column of index \(i\) . We denote this square matrix by \(X^{ji}\) . Obviously \(\rho_{\mathrm{H},\mathrm{H}'} = (-1)^{i-1}\) and \(\rho_{\mathrm{R},\mathrm{R}'} = (-1)^{j-1}\) ; therefore (21) becomes in this case

\[
\det X = \sum_ {j = 1} ^ {n} (- 1) ^ {i + j} \xi_ {j i} \det (X ^ {j i})\tag{23}
\]

and we obtain similarly from (22)

\[
\sum_ {j = 1} ^ {n} (- 1) ^ {j} \xi_ {j i} \det (X ^ {j k}) = 0 \quad \text { for } k \neq i.\tag{24}
\]

Formula (23) is known as the expansion of the determinant of X by the column of index i. The scalar \((-1)^{i+j}\det(X^{ji})\) is called the cofactor of indices j and i (or, by an abuse of language, the cofactor of \(\xi_{ji}\) ) in X.

The matrix of cofactors of X is the matrix

\[
Y = ((- 1) ^ {i + j} \det (X ^ {j i}))\tag{25}
\]

whose element in the j-th row and i-th column is the cofactor of indices j and i. Formulae (23) and (24) are equivalent to the formula

\[
{ } ^ { t } Y . X = ( \operatorname * { d e t } X ) I _ { n } .\tag{26}
\]

Therefore:

PROPOSITION 12. For every invertible square matrix \(X\) of type \((n, n)\) , the inverse of \(X\) is given by the formula

\[
X ^ {- 1} = (\det X) ^ {- 1}. ^ {t} Y\tag{27}
\]

where Y is the cofactor matrix of X.

By considering the transpose of X and using Proposition 8 of no. 5, the Laplace expansions could be obtained relative to two complementary sets of rows and, in particular, the expansion of det X by a row, thus there are formulae equivalent to

\[
X. ^ {t} Y = (\det X) I _ {n},\tag{28}
\]

in the above notation.

It is easily verified that if X is the matrix of an endomorphism u of a free A-module M of dimension n with respect to a basis \((e_{i})_{1\leq i\leq n}\) , \(^{t}Y\) is the matrix of the endomorphism \(\tilde{u}\) of M defined by the following condition: for every set of n elements \(x, y_{2}, \ldots, y_{n}\) of M,

\[
\tilde {u} (x) \wedge y _ {2} \wedge \dots \wedge y _ {n} = x \wedge u (y _ {2}) \wedge \dots \wedge u (y _ {n}).
\]

\(\tilde{u}\) is called the cotranspose of u (cf.~§ 11, no. 11, Corollary to Proposition 13).

Examples. (1) Vandermonde determinant. Given a sequence \((\zeta_{i})_{1\leqslant i\leqslant n}\) of \(n\) elements of A, the Vandermonde determinant of this sequence is the determinant

\[
\mathrm{V} (\zeta_ {1}, \zeta_ {2}, \dots , \zeta_ {n}) = \left| \begin{array}{c c c c} 1 & 1 & \dots & 1 \\ \zeta_ {1} & \zeta_ {2} & \dots & \zeta_ {n} \\ \zeta_ {1} ^ {2} & \zeta_ {2} ^ {2} & \dots & \zeta_ {n} ^ {2} \\ \cdot & \cdot & \cdot & \cdot \\ \zeta_ {1} ^ {n - 1} & \zeta_ {2} ^ {n - 1} & \dots & \zeta_ {n} ^ {n - 1} \end{array} \right|
\]

We shall show that

\[
\mathrm{V} \left(\zeta_ {1}, \zeta_ {2}, \dots , \zeta_ {n}\right) = \prod_ {i <   j} \left(\zeta_ {j} - \zeta_ {i}\right).\tag{29}
\]

Since the proposition is immediate for \(n = 1\) , we argue by induction on \(n\) .

For each index \(k \geqslant 2\) , we subtract from the row of index k the row of index k - 1 multiplied by \(\zeta_{1}\) ; the value of the determinant is unaltered and hence

\[
\mathrm{V} \left(\zeta_ {1}, \zeta_ {2}, \dots , \zeta_ {n}\right) = \left| \begin{array}{c c c c c c c c c} 1 & 1 & & \dots & 1 \\ 0 & \zeta_ {2} - \zeta_ {1} & & \dots & \zeta_ {n} - \zeta_ {1} \\ 0 & \zeta_ {2} \left(\zeta_ {2} - \zeta_ {1}\right) & & \dots & \zeta_ {n} \left(\zeta_ {n} - \zeta_ {1}\right) \\ . & . & . & . & . & . & . & . & . \\ 0 & \zeta_ {2} ^ {n - 2} \left(\zeta_ {2} - \zeta_ {1}\right) & & \dots & \zeta_ {n} ^ {n - 2} \left(\zeta_ {n} - \zeta_ {1}\right) \end{array} \right|
\]

whence, expanding by the first column and then taking out the factor \(\zeta_{k}-\zeta_{1}\) from the column of index k-1 from the minor thus obtained ( \(2\leqslant k\leqslant n\) )

\[
\mathrm{V} \left(\zeta_ {1}, \dots , \zeta_ {n}\right) = \left(\zeta_ {2} - \zeta_ {1}\right) \left(\zeta_ {3} - \zeta_ {1}\right) \dots \left(\zeta_ {n} - \zeta_ {1}\right) V \left(\zeta_ {2}, \dots , \zeta_ {n}\right)
\]

which establishes (29) by induction.

\begin{enumerate}
\def\labelenumi{(\arabic{enumi})}
\setcounter{enumi}{1}
\tightlist
\item
  Consider a square matrix of order n which is presented in the form of an ``upper triangular matrix of matrices'' (II, § 10, no. 7, Example IV)
\end{enumerate}

\[
X = \left( \begin{array}{c c} Y & T \\ 0 & Z \end{array} \right)
\]

We show that

\begin{enumerate}
\def\labelenumi{(\arabic{enumi})}
\setcounter{enumi}{29}
\tightlist
\item
\end{enumerate}

\[
\det X = (\det Y) (\det Z).
\]

Let \(n\) be the order of the matrix \(X, h\) that of \(Y, (e_i)_{1 \leqslant i \leqslant n}\) the canonical basis of \(A^n\) and \(x_i\) ( \(1 \leqslant i \leqslant n\) ) the columns of \(X\) ; the hypothesis implies that the columns \(x_1, \ldots, x_h\) belong to the submodule of \(A^n\) with basis \(e_1, \ldots, e_h\) and then by definition (no. 3, formula (6))

\[
x _ {1} \wedge x _ {2} \wedge \dots \wedge x _ {h} = (\det Y) e _ {1} \wedge e _ {2} \wedge \dots \wedge e _ {h}.
\]

On the other hand, for every index i \textgreater{} h, we can write \(x_{i} = y_{i} + z_{i}\) , where \(y_{i}\) is a linear combination of \(e_{1}, e_{2}, \ldots, e_{h}\) and \(z_{i}\) is a linear combination of \(e_{h+1}, \ldots, e_{n}\) . By (30), \(x_{1} \wedge x_{2} \wedge \cdots \wedge x_{h} \wedge y_{i} = 0\) for all i \textgreater{} h, therefore

\[
x _ {1} \wedge x _ {2} \wedge \dots \wedge x _ {n} = (\det Y) e _ {1} \wedge e _ {2} \wedge \dots \wedge e _ {h} \wedge z _ {h + 1} \wedge \dots \wedge z _ {n}.
\]

But by definition

\[
z _ {h + 1} \wedge z _ {h + 2} \wedge \dots \wedge z _ {n} = (\det Z) e _ {h + 1} \wedge e _ {h + 2} \wedge \dots \wedge e _ {n}
\]

whence formula (30).

By induction on \(p\) , it follows that if \(X\) is of the form of an upper triangular matrix of matrices:

\[
X = \left( \begin{array}{c c c c c} X _ {1 1} & X _ {1 2} & \ldots & X _ {1 p} \\ 0 & X _ {2 2} & \ldots & X _ {2 p} \\ \cdot & \cdot & \cdot & \cdot & \cdot \\ 0 & 0 & \ldots & X _ {p p} \end{array} \right)\tag{31}
\]

\[
\det X = (\det X _ {1 1}) (\det X _ {2 2}) \dots (\det X _ {p p}).
\]

This can be applied in particular to a triangular matrix (where all the \(X_{ii}\) are of order 1) and more particularly to a diagonal matrix:

\[
\det (\operatorname{diag} \left(\alpha_ {1}, \alpha_ {2}, \dots , \alpha_ {n}\right)) = \alpha_ {1} \alpha_ {2} \dots \alpha_ {n}.\tag{32}
\]

\begin{enumerate}
\def\labelenumi{(\arabic{enumi})}
\setcounter{enumi}{2}
\tightlist
\item
  Let M, M' be two free A-modules of respective dimensions n, \(n'\) , u an endomorphism of M and \(u'\) an endomorphism of M'. Then
\end{enumerate}

\[
\det (u \otimes u ^ {\prime}) = (\det u) ^ {n ^ {\prime}} (\det u ^ {\prime}) ^ {n}.\tag{33}
\]

For we can write \(u \otimes u' = (u \otimes 1_{\mathbf{M}'}) \circ (1_{\mathbf{M}} \otimes u')\) and are then led to the case where one of the two endomorphisms \(u, u'\) is the identity. For example if \(u' = 1_{\mathbf{M}'}\) and \(X\) is the matrix of \(u\) with respect to a basis \((e_i)\) of \(\mathbf{M}\) , then the matrix of \(u \otimes 1_{\mathbf{M}'}\) with respect to the tensor product of \((e_i)\) and a basis of \(\mathbf{M}'\) can be written as a matrix (with \(n'\) rows and \(n'\) columns) of matrices with \(n\) rows and \(n\) columns

\[
\left( \begin{array}{c c c c} X & 0 & \ldots & 0 \\ 0 & X & \ldots & 0 \\ \cdot & \cdot & \cdot & \cdot \\ 0 & 0 & \ldots & X \end{array} \right)
\]

whence, by virtue of Example 2,

\[
\det (u \otimes 1 _ {\mathbf {M} ^ {\prime}}) = (\det X) ^ {n ^ {\prime}} = (\det u) ^ {n ^ {\prime}}
\]

which immediately gives formula (33).

